% ============================================================
\chapter{Attracteurs \'Etranges}
\label{ch:attracteurs}
% ============================================================

\section{Introduction}

En 1971, David Ruelle et Floris Takens proposent le terme d'\og attracteur \'etrange \fg{} pour d\'ecrire les ensembles g\'eom\'etriques qui gouvernent la dynamique des syst\`emes chaotiques. L'attracteur de Lorenz, d\'ecouvert en 1963, en est le prototype~: un ensemble en forme de papillon, de dimension fractale $\approx 2{,}06$, sur lequel les trajectoires s'enroulent sans jamais se recouper. L'attracteur de H\'enon, l'attracteur de R\"ossler, et bien d'autres suivront. Ce qui rend ces objets \og \'etranges \fg, c'est la coexistence de deux comportements contradictoires~: l'attraction (les orbites voisines convergent vers l'attracteur) et le chaos (une fois sur l'attracteur, les orbites divergent exponentiellement). G\'eom\'etriquement, cette coexistence se traduit par une structure fractale auto-similaire, r\'esultat de l'\'etirement et du repliement r\'ep\'et\'es de l'espace des phases.

\begin{intuition}
Un attracteur \'etrange fonctionne comme une machine \`a p\'etrir la p\^ate :
il \'etire les orbites voisines (cr\'eant la sensibilit\'e aux conditions
initiales) puis les replie dans un volume born\'e (assurant la compacit\'e).
Le r\'esultat est une structure feuillet\'ee infiniment fine --- un fractal.
\end{intuition}

\section{D\'efinitions formelles}

\begin{definition}[Attracteur]
\index{attracteur}
Soit $\varphi_t$ un flot sur $\R^n$. Un ensemble compact invariant $A$ est
un \textbf{attracteur} s'il existe un voisinage ouvert $U \supset A$ (bassin
d'attraction) tel que :
\begin{enumerate}
    \item $\varphi_t(U) \subset U$ pour tout $t > 0$ ;
    \item $A = \bigcap_{t > 0} \varphi_t(U)$ ;
    \item $A$ est topologiquement transitif (contient une orbite dense).
\end{enumerate}
\end{definition}

\begin{definition}[Attracteur \'etrange]
Un attracteur est dit \textbf{\'etrange}\index{attracteur \'etrange!d\'efinition}
s'il poss\`ede une structure fractale (dimension non enti\`ere) et s'il
exhibe une sensibilit\'e aux conditions initiales (au moins un exposant
de Lyapunov strictement positif).
\end{definition}

\begin{attention}[title=\textbf{\'Etrange $\neq$ chaotique}]
Il existe des attracteurs \'etranges non chaotiques (dimension fractale mais
exposants de Lyapunov n\'egatifs ou nuls) et des attracteurs chaotiques non
\'etranges (chaos sur un tore lisse). Les deux notions sont logiquement
ind\'ependantes.
\end{attention}

\section{L'attracteur de Lorenz}

\begin{definition}[Syst\`eme de Lorenz]
\index{syst\`eme de Lorenz}
Le \textbf{syst\`eme de Lorenz} (1963) est
\[
    \begin{cases}
    \dot{x} = \sigma(y - x), \\
    \dot{y} = \rho x - y - xz, \\
    \dot{z} = xy - \beta z,
    \end{cases}
\]
avec les param\`etres classiques $\sigma = 10$, $\rho = 28$, $\beta = 8/3$.
\end{definition}

\begin{theoreme}[Tucker, 2002]
\index{th\'eor\`eme de Tucker}
Pour les param\`etres classiques $(\sigma, \rho, \beta) = (10, 28, 8/3)$, le
syst\`eme de Lorenz poss\`ede un \textbf{attracteur \'etrange robuste}. Ce
r\'esultat, d\'emontr\'e par assistance informatique rigoureuse, a confirm\'e
une conjecture ouverte depuis 1963.
\end{theoreme}

\begin{proposition}[Propri\'et\'es de l'attracteur de Lorenz]
L'attracteur de Lorenz satisfait :
\begin{enumerate}
    \item \textbf{Dissipativit\'e} : $\nabla \cdot f = -(\sigma + 1 + \beta) < 0$,
    donc les volumes se contractent exponentiellement.
    \item \textbf{Sym\'etrie} : invariance par $(x, y, z) \mapsto (-x, -y, z)$.
    \item \textbf{Spectre de Lyapunov} : $\lambda_1 \approx 0{,}91$,
    $\lambda_2 = 0$, $\lambda_3 \approx -14{,}57$.
    \item \textbf{Dimension fractale} : $D_{KY} \approx 2{,}06$.
\end{enumerate}
\end{proposition}

\begin{formules}[title=\textbf{Points d'\'equilibre du syst\`eme de Lorenz}]
\begin{align*}
    &C_0 = (0, 0, 0), \quad \text{instable pour } \rho > 1. \\
    &C_{\pm} = (\pm\sqrt{\beta(\rho - 1)}, \pm\sqrt{\beta(\rho - 1)}, \rho - 1),
    \quad \text{instables pour } \rho > \rho_H = \sigma\frac{\sigma + \beta + 3}{\sigma - \beta - 1}.
\end{align*}
\end{formules}

\section{Le syst\`eme de R\"ossler}

\begin{definition}[Syst\`eme de R\"ossler]
\index{syst\`eme de R\"ossler}
Le \textbf{syst\`eme de R\"ossler} (1976) est
\[
    \begin{cases}
    \dot{x} = -y - z, \\
    \dot{y} = x + ay, \\
    \dot{z} = b + z(x - c),
    \end{cases}
\]
avec les param\`etres typiques $a = 0{,}2$, $b = 0{,}2$, $c = 5{,}7$.
\end{definition}

\begin{remarque}
Contrairement \`a l'attracteur de Lorenz qui a deux lobes, l'attracteur de
R\"ossler a une structure en bande unique avec un m\'ecanisme de repliement
plus simple. Il est souvent utilis\'e comme mod\`ele minimal du chaos \`a
bande unique.
\end{remarque}

\begin{proposition}
Pour les param\`etres standards, le syst\`eme de R\"ossler exhibe une cascade
de doublements de p\'eriode lorsque $c$ augmente : orbite p\'eriodique
($c \approx 2{,}5$), doublement ($c \approx 3{,}5$), chaos ($c \approx 4{,}2$),
fen\^etre p\'eriodique ($c \approx 5{,}2$), chaos d\'evelopp\'e ($c \approx 5{,}7$).
\end{proposition}

\section{L'application de H\'enon}

\begin{definition}[Application de H\'enon]
\index{application de H\'enon}
L'\textbf{application de H\'enon} est le diff\'eomorphisme du plan
\[
    H(x, y) = (1 - ax^2 + y, \; bx),
\]
avec les param\`etres classiques $a = 1{,}4$, $b = 0{,}3$.
\end{definition}

\begin{proposition}[Propri\'et\'es de H\'enon]
\begin{enumerate}
    \item Le jacobien est constant : $\det DH = -b$. L'application contracte
    les aires d'un facteur $\abs{b} = 0{,}3$.
    \item L'attracteur a une dimension fractale $D \approx 1{,}26$.
    \item Les exposants de Lyapunov sont $\lambda_1 \approx 0{,}42$ et
    $\lambda_2 \approx -1{,}62$.
    \item L'attracteur poss\`ede une structure de Cantor dans la direction
    transverse (produit d'une courbe par un ensemble de Cantor).
\end{enumerate}
\end{proposition}

\begin{theoreme}[Benedicks-Carleson, 1991]
\index{th\'eor\`eme de Benedicks-Carleson}
Pour un ensemble de param\`etres $(a, b)$ de mesure de Lebesgue positive avec
$b$ petit, l'application de H\'enon poss\`ede un attracteur \'etrange portant
une unique mesure SRB (Sinai-Ruelle-Bowen) ergodique.
\end{theoreme}

\section{Dimension fractale}

\begin{definition}[Dimension de bo\^ite (box-counting)]
\index{dimension!box-counting}
Soit $A \subset \R^n$ un ensemble born\'e. La \textbf{dimension de bo\^ite} est
\[
    D_0 = \lim_{\epsilon \to 0} \frac{\ln N(\epsilon)}{-\ln \epsilon},
\]
o\`u $N(\epsilon)$ est le nombre minimal de bo\^ites de c\^ot\'e $\epsilon$
n\'ecessaires pour couvrir $A$.
\end{definition}

\begin{definition}[Dimension de Hausdorff]
\index{dimension!Hausdorff}
La \textbf{dimension de Hausdorff} de $A$ est
\[
    D_H = \inf\left\{ d \geq 0 : \mathcal{H}^d(A) = 0 \right\},
\]
o\`u $\mathcal{H}^d(A) = \lim_{\epsilon \to 0} \inf \left\{
\sum_{i} (\mathrm{diam}\, U_i)^d : A \subset \bigcup_i U_i, \;
\mathrm{diam}\, U_i < \epsilon \right\}$ est la mesure de Hausdorff
$d$-dimensionnelle.
\end{definition}

\begin{proposition}[In\'egalit\'e entre dimensions]
Pour tout ensemble born\'e $A \subset \R^n$ :
\[
    D_H(A) \leq \underline{D}_0(A) \leq \overline{D}_0(A),
\]
o\`u $\underline{D}_0$ et $\overline{D}_0$ sont les dimensions de bo\^ite
inf\'erieure et sup\'erieure.
\end{proposition}

\begin{exemple}[Ensemble de Cantor triadique]
L'ensemble de Cantor triadique $C$ a pour dimension
$D_H = D_0 = \frac{\ln 2}{\ln 3} \approx 0{,}631$.
En effet, \`a l'\'etape $n$, on a $N(\epsilon) = 2^n$ intervalles de longueur
$\epsilon = 3^{-n}$, donc $D_0 = \lim_{n \to \infty} \frac{n \ln 2}{n \ln 3}$.
\end{exemple}

\begin{definition}[Dimension de corr\'elation]
\index{dimension!corr\'elation}
La \textbf{dimension de corr\'elation} est d\'efinie par
\[
    D_2 = \lim_{\epsilon \to 0} \frac{\ln C(\epsilon)}{\ln \epsilon},
\]
o\`u $C(\epsilon) = \lim_{N \to \infty} \frac{1}{N^2} \#\{(i, j) :
\norm{x_i - x_j} < \epsilon\}$ est l'int\'egrale de corr\'elation
(Grassberger-Procaccia).
\end{definition}

\section{Attracteurs dans les syst\`emes physiques}

\begin{exemple}[Convection de Rayleigh-B\'enard]
Le syst\`eme de Lorenz est une troncation des \'equations de
Navier-Stokes pour la convection de Rayleigh-B\'enard entre deux plaques
horizontales chauff\'ees. Le param\`etre $\rho$ correspond au nombre de
Rayleigh normalis\'e.
\end{exemple}

\begin{exemple}[Circuit de Chua]
\index{circuit de Chua}
Le \textbf{circuit de Chua} est un circuit \'electronique simple
(condensateurs, inductance, r\'esistance non lin\'eaire) qui produit un
attracteur \'etrange \`a double enroulement. C'est le premier syst\`eme
physique pour lequel le chaos a \'et\'e rigoureusement d\'emontr\'e.
\end{exemple}

\begin{codeblock}[title=\textbf{Visualisation de l'attracteur de Lorenz}]
\begin{minted}{python}
import numpy as np
from scipy.integrate import solve_ivp
import matplotlib.pyplot as plt

def lorenz(t, state, sigma=10, rho=28, beta=8/3):
    x, y, z = state
    return [sigma*(y - x), rho*x - y - x*z, x*y - beta*z]

sol = solve_ivp(lorenz, [0, 100], [1, 1, 1],
                max_step=0.01, t_eval=np.linspace(0, 100, 100000))
fig = plt.figure(figsize=(10, 8))
ax = fig.add_subplot(111, projection='3d')
ax.plot(sol.y[0], sol.y[1], sol.y[2], lw=0.3, color='darkblue')
ax.set_xlabel('$x$'); ax.set_ylabel('$y$'); ax.set_zlabel('$z$')
ax.set_title('Attracteur de Lorenz')
plt.savefig("lorenz_attractor.pdf")
plt.show()
\end{minted}
\end{codeblock}

\section{Exercices}

\begin{exercice}[Dissipativit\'e]
Montrer que le syst\`eme de Lorenz est dissipatif et que tout ensemble de
volume initial $V_0$ \'evolue en un volume $V(t) = V_0 e^{-(\sigma + 1 + \beta)t}$.
En d\'eduire l'existence d'un attracteur de mesure de Lebesgue nulle.
\end{exercice}

\begin{exercice}[Dimension de bo\^ite]
Calculer la dimension de bo\^ite de l'ensemble $\{0\} \cup \{1/n : n \in \N^*\}$
et de l'\'eponge de Menger.
\end{exercice}

\begin{exercice}[Attracteur de H\'enon]
\'Ecrire un programme qui trace l'attracteur de H\'enon pour $a = 1{,}4$,
$b = 0{,}3$ en it\'erant $10^6$ fois. Estimer sa dimension de corr\'elation
par la m\'ethode de Grassberger-Procaccia.
\end{exercice}

\begin{exercice}[Section de Poincar\'e]
Pour le syst\`eme de R\"ossler, tracer la section de Poincar\'e dans le plan
$z = z_{\max}/2$ et montrer qu'elle a l'apparence d'une application
unidimensionnelle (presque).
\end{exercice}

\begin{exercice}[Stabilit\'e structurelle]
L'attracteur de Lorenz est-il structurellement stable ? Discuter en comparant
avec le fer \`a cheval de Smale et justifier la notion d'attracteur robuste
(au sens de Morales-Paci\'fico-Pujals).
\end{exercice}
